% Part:sets-functions-relations
% Chapter: size-of-sets
% Section: enumerations

\documentclass[../../../include/open-logic-section]{subfiles}

\begin{document}

\olfileid{sfr}{siz}{enm}

\olsection{可数集与枚举}

\begin{editorial}
  本节讨论集合的枚举，将其定义为以 $\PosInt$ 为定义域的满射。讲解较为缓慢，适合数学背景较浅的读者。在 \olref[enm-alt]{sec} 中给出了一个更简洁的替代版本，该版本以不同方式定义了枚举：即与 $\Nat$（或其初始段）的双射。
\end{editorial}

\begin{explain}
我们已经通过列出元素的方式给出过一些集合的例子。让我们更一般地讨论，即便集合是无限的，我们如何以及何时能够列出它的元素。
\end{explain}

\begin{defn}[枚举，非正式定义]
非正式地，集合~$A$ 的一个\emph{枚举}是一个由~$A$ 的元素构成的（可能无限的）列表，使得 $A$ 的每一个元素都在该列表的某个有限位置出现。如果 $A$ 有一个枚举，则称 $A$ 是\emph{可数的}。
\end{defn}

\begin{explain}
关于枚举的几点说明：
\begin{enumerate}
\item 我们只将有开头、且除第一个元素外每个元素都恰好有一个直接前趋的列表视为枚举。换言之，在列表的第一个元素与任何其他元素之间，只有有限多个元素。特别地，这意味着枚举中的每个元素都有一个有限的位置：第一个元素占据位置~$1$，第二个占据位置~$2$，依此类推。
\item 对于同一个集合~$A$，我们可以有不同的枚举，它们按元素出现的顺序不同而不同：$4$, $1$, $25$, $16$,~$9$ 和 $1$, $4$, $9$, $16$,~$25$ 一样，都枚举了前五个平方数（的集合）。
\item 冗余的枚举仍然是枚举：$1$, $1$, $2$, $2$, $3$, $3$,~\dots{} 所枚举的集合与 $1$, $2$, $3$,~\dots{} 所枚举的相同。
\item 在具体指定一个枚举时，顺序和冗余\emph{确实}有影响：我们可以从 $1$, $2$, $3$, $1$, \dots{} 开始枚举正整数，但当按照标准方式 $1$, $2$, $3$, $4$,~\dots{} 枚举时，其模式更容易看清。
\item 枚举必须有开始：\dots, $3$, $2$, $1$ 不是正整数的枚举，因为它没有第一个元素。要理解这如何从非正式定义得出，不妨问问自己：“数字 76 在这个列表的什么位置出现？”
\item 下面这个不是正整数的枚举：$1$, $3$, $5$, \dots, $2$, $4$, $6$, \dots\@ 问题在于，偶数出现在位置 $\infty + 1$、$\infty + 2$、$\infty + 3$，而不是有限的位置上。
\item 空集是可数的：空列表便枚举了它！
\end{enumerate}
\end{explain}

\begin{prop}
如果 $A$ 有一个枚举，那么它有一个无重复的枚举。
\end{prop}

\begin{proof}
假设 $A$ 有一个枚举 $x_1$, $x_2$, \dots{}，其中每个 $x_i$ 都是~$A$ 的一个元素。我们可以通过移除重复元素来得到一个没有重复的枚举。例如，我们可以将其转化为一个新的枚举：如果 $x_i$ 是不在 $x_1$, \dots, $x_{i-1}$ 之中的 $A$ 的元素，我们就将其列入；如果它已经出现在 $x_1$, \dots,~$x_{i-1}$ 之中，我们便将其跳过。
\end{proof}

上述论证表明，为了能很好地处理枚举和可数集并证明关于它们的性质，我们需要一个更精确的定义。下面给出这个定义。

\begin{defn}[枚举，形式化定义]
集合 $A \neq \emptyset$ 的一个\emph{枚举}是任意满射函数 $f \colon \PosInt \to A$。
\end{defn}

\begin{explain}
让我们来确认，上述形式定义与使用（可能无限的）列表的非形式定义是等价的。首先，任何从 $\PosInt$ 到集合~$A$ 的满射函数都枚举了~$A$。这样的函数决定了上文非形式定义的枚举：列表 $f(1)$, $f(2)$, $f(3)$, \dots。由于 $f$ 是满射的，$A$ 的每个元素必定是某个 $n \in \PosInt$ 对应的 $f(n)$ 的值。因此，$A$ 的每个元素都会在该列表的某个有限位置出现。由于该函数可能不是单射的，列表可能有重复，但这是可以接受的（如上所述）。

另一方面，给定一个枚举了~$A$ 所有元素的列表，我们可以定义一个满射函数 $f\colon \PosInt \to A$，令 $f(n)$ 为列表的第 $n$ 个元素（如果列表没有第 $n$ 个元素，则令其为列表的最后一个元素）。唯一无法通过此法产生满射函数的情况是 $A$ 为空集（从而列表也为空）。因此，每个非空列表都决定了一个满射函数 $f\colon \PosInt \to A$。
\end{explain}

\begin{defn}
  \ollabel{defn:enumerable}
  集合~$A$ 是\emph{可数的}，当且仅当 $A$ 是空集或者存在一个 $A$ 的枚举。
\end{defn}

\begin{ex}
枚举正整数（$\PosInt$）的函数很简单，就是由 $f(n) = n$ 给出的恒等函数。枚举自然数 $\Nat$ 的函数是 $g(n) = n - 1$。
\end{ex}

\begin{ex}
由
\begin{align*}
f(n) & = 2n \text{ 与}\\
g(n) & = 2n - 1
\end{align*}
给出的函数 $f\colon \PosInt \to \PosInt$ 和 $g \colon \PosInt \to
\PosInt$ 分别枚举了正偶整数和正奇整数。然而，这两个函数都不是 $\PosInt$ 的枚举，因为它们都不是满射的。
\end{ex}

\begin{prob}
定义正整数的平方 $1$, $4$, $9$, $16$, \dots 所构成的集合的枚举。
\end{prob}

\begin{ex}
函数 $f(n) = (-1)^{n} \lceil \frac{(n-1)}{2}\rceil$（其中
$\lceil x \rceil$ 表示\emph{上取整}函数，即将
$x$ 向上舍入到最近的整数）枚举了整数集~$\Int$。注意 $f$ 是如何通过在正整数和负整数之间“来回跳跃”来生成 $\Int$ 的值的：
\[
\begin{array}{c c c c c c c c}
f(1) & f(2) & f(3) & f(4) & f(5) & f(6) & f(7) & \dots \\ \\
- \lceil \tfrac{0}{2} \rceil & \lceil \tfrac{1}{2}\rceil & - \lceil \tfrac{2}{2} \rceil & \lceil \tfrac{3}{2} \rceil & - \lceil \tfrac{4}{2} \rceil  & \lceil \tfrac{5}{2}
\rceil & - \lceil \tfrac{6}{2} \rceil & \dots \\ \\
0 & 1 & -1 & 2 & -2 & 3 & \dots
\end{array}
\]
你也可以将 $f$ 看作是如下分段定义的：
\[
f(n) = \begin{cases}
  0 & \text{若 $n = 1$}\\
  n/2 & \text{若 $n$ 为偶数}\\
  -(n-1)/2 & \text{若 $n$ 为奇数且 $>1$}
  \end{cases}
\]
\end{ex}

\begin{prob}
  证明：如果 $A$ 和 $B$ 是可数的，那么 $A \cup B$ 也是可数的。为此，假设存在满射函数 $f\colon \PosInt \to
  A$ 和 $g\colon \PosInt \to B$，然后定义一个满射
  函数~$h\colon \PosInt \to A \cup B$ 并证明它是满射的。同时考虑 $A$ 或~$B = \emptyset$ 的情况。
\end{prob}

\begin{prob}
  证明：如果 $B \subseteq A$ 且 $A$ 是可数的，那么~$B$ 也是可数的。为此，假设存在一个满射函数 $f\colon \PosInt \to
  A$。定义一个满射函数~$g\colon \PosInt \to B$ 并证明它是满射的。如果 $B = \emptyset$，情况如何？
\end{prob}

\begin{prob}
  通过对 $n$ 进行归纳，证明：如果 $A_1$, $A_2$, \dots, $A_n$ 都是可数的，那么 $A_1 \cup \dots \cup A_n$ 也是可数的。你可以使用如下事实：如果两个集合 $A$ 和~$B$ 是可数的，那么~$A \cup
  B$ 也是可数的。
\end{prob}

虽然列举集合元素时从第 $1$ 个元素开始计数可能更自然，但数学家们喜欢用自然数~$\Nat$ 来计数。他们谈论列表的第 $0$ 个、第 $1$ 个、第 $2$ 个等元素。相应地，我们可以将枚举定义为一个从 $\Nat$ 到~$A$ 的满射函数。当然，这两种定义是等价的。

\begin{prop}\ollabel{prop:enum-shift}
  存在满射 $f\colon \PosInt \to A$ 当且仅当存在满射 $g\colon \Nat \to A$。
\end{prop}

\begin{proof}
  给定一个满射 $f\colon \PosInt \to A$，我们可以对所有的 $n \in \Nat$ 定义 $g(n) =
  f(n+1)$。容易看出 $g\colon \Nat
  \to A$ 是满射的。反之，给定一个满射 $g\colon
  \Nat \to A$，定义 $f(n) = g(n-1)$。
\end{proof}

由此我们得到以下推论：

\begin{cor}\ollabel{cor:enum-nat}
集合 $A$ 是可数的，当且仅当 $A$ 是空集或者存在一个满射函数 $f\colon \Nat \to A$。
\end{cor}

我们上面讨论过，一个集合~$A$ 的元素列表可以转化为一个无重复列表。枚举也是如此，只是要严格地加以形式化与证明会稍微困难一些。任何函数 $f\colon \PosInt \to A$ 都必须在所有 $n \in \PosInt$ 上有定义。若~$A$ 中只有有限多个元素，则我们显然不可能有一个定义在~$\PosInt$ 的无穷多个元素上、以~$A$ 的所有元素为值且绝不取同一个值两次的函数。在这种情况下，即无重复列表为有限时，我们必须为~$f$ 选择一个不同的定义域，一个仅含有限多个元素的定义域。无重复意味着 $f$ 必须是单射的。既然它同时也是满射的，我们实际上是在寻找某个有限集 $\{1, \dots, n\}$ 或 $\PosInt$ 与~$A$ 之间的一个双射。

\begin{prop}\ollabel{prop:enum-bij}
若 $f\colon \PosInt \to A$ 是满射（即~$A$ 的一个枚举），则存在一个双射 $g\colon Z \to A$，其中 $Z$ 是~$\PosInt$ 或对于某个~$n \in \PosInt$ 为 $\{1, \dots, n\}$。
\end{prop}

\begin{proof}
  我们递归地定义函数 $g$：令 $g(1) = f(1)$。若 $g(i)$ 已有定义，则令 $g(i+1)$ 为 $f(1)$, $f(2)$, \dots{} 中第一个不在 $g(1)$, \dots, $g(i)$ 之中的值（若存在这样的值）。若 $A$ 刚好有 $n$ 个元素，则 $g(1)$, \dots, $g(n)$ 就都有定义，于是我们就定义了一个函数 $g\colon \{1, \dots, n\}
  \to A$。若 $A$ 有无穷多个元素，则对于任意 $i$，在枚举 $f(1)$, $f(2)$, \dots 中必定有一个尚未出现在 $g(1)$, \dots, $g(i)$ 中的~$A$ 的元素。在这种情况下，我们就定义了一个函数 $g\colon \PosInt \to A$。
  
  函数 $g$ 是满射的，因为~$A$ 的任何元素都在 $f(1)$, $f(2)$, \dots{} 之中（由于 $f$ 是满射），因此最终总会成为某个 $g(i)$ 的值。它也是单射的，因为若有 $j < i$ 使得 $g(j) = g(i)$，则 $g(i)$ 就已经出现在 $g(1)$, \dots, $g(i-1)$ 之中，这与我们定义~$g$ 的方式相矛盾。
\end{proof}

\begin{cor}\ollabel{cor:enum-nat-bij}
集合 $A$ 是可数的，当且仅当 $A$ 是空集，或存在一个双射 $f\colon N \to A$，其中 $N = \Nat$ 或对于某个 $n \in \Nat$ 为 $N = \{0, \dots, n\}$。
\end{cor}

\begin{proof}
$A$ 是可数的，当且仅当 $A$ 是空集，或存在一个满射 $f\colon \PosInt \to A$。根据 \olref{prop:enum-bij}，后者成立当且仅当存在一个双射函数~$f\colon Z \to A$，其中 $Z = \PosInt$ 或对于某个 $n \in \PosInt$ 为 $Z = \{1, \dots, n\}$。根据与 \olref{prop:enum-shift} 的证明中相同的论证，这又等价于存在一个双射 $g\colon N \to A$，其中 $N = \Nat$ 或 $N = \{0, \dots, n-1\}$。
\end{proof}

\begin{prob}
  根据 \olref[sfr][siz][enm]{defn:enumerable}，集合 $A$ 是可数的，当且仅当 $A = \emptyset$ 或存在一个满射 $f\colon \PosInt \to A$。也可以用如下方式精确定义“可数集”：一个集合是可数的，当且仅当存在单射函数 $g\colon A \to \PosInt$。请证明这两个定义是等价的，即证明：存在单射函数 $g\colon A \to \PosInt$，当且仅当 $A = \emptyset$ 或存在满射 $f\colon \PosInt \to A$。
\end{prob}